% !TeX encoding = UTF-8
\documentclass[variant=guia,cover=institutional,mode=screen]{ua-engineering-v6-article}
% !TeX encoding = UTF-8
\UASetUniversity{Universidad Austral}
\UASetFaculty{Facultad de Ingeniería}
\UASetDepartment{Departamento de Computación}
\UASetCampus{Pilar}
\UASetCareer{Ingeniería Informática}
\UASetCommission{Comisión A}
\UASetProfessor{Equipo docente}
\UASetSemester{Segundo cuatrimestre 2026}
\UASetCourse{Sistemas Distribuidos}
\UASetDate{2026}


\UASetClassNumber{03}
\UASetClassTitle{Consistencia y replicación}
\UASetDocKind{Guía resuelta}

\title{Clase 03 --- Guía resuelta}
\author{\UAProfessor}
\date{\UADate}

\begin{document}

\section{Criterio de lectura}
Las resoluciones son modelos de razonamiento. Deben verificarse contra los supuestos del caso y no memorizarse como recetas independientes del dominio.

\section{Ejercicios y resoluciones completas}
\begin{uaexercise}
Definir replicación y explicar sus motivaciones principales.
\end{uaexercise}
\begin{uaanswer}
La idempotencia separa intentos físicos de una operación lógica. Se usa una clave estable, se persiste el resultado junto con el efecto y los reintentos devuelven ese mismo resultado. El costo es estado adicional, política de expiración y coordinación con el almacenamiento.

La evidencia apropiada sería historia de lecturas/escrituras, staleness y comparación de quórums. También debe indicarse una condición que refute la elección o haga preferible una solución más simple.
\end{uaanswer}

\begin{uaexercise}
Explicar por qué la divergencia es inevitable en un sistema asíncrono.
\end{uaexercise}
\begin{uaanswer}
La idempotencia separa intentos físicos de una operación lógica. Se usa una clave estable, se persiste el resultado junto con el efecto y los reintentos devuelven ese mismo resultado. El costo es estado adicional, política de expiración y coordinación con el almacenamiento.

El argumento debe identificar la hipótesis, construir la cadena lógica o el contraejemplo y concluir exactamente qué propiedad queda demostrada. Una intuición sin secuencia verificable no es suficiente.

La evidencia apropiada sería historia de lecturas/escrituras, staleness y comparación de quórums. También debe indicarse una condición que refute la elección o haga preferible una solución más simple.
\end{uaanswer}

\begin{uaexercise}
Construir un ejemplo con dos réplicas que divergen tras escrituras concurrentes.
\end{uaexercise}
\begin{uaanswer}
La idempotencia separa intentos físicos de una operación lógica. Se usa una clave estable, se persiste el resultado junto con el efecto y los reintentos devuelven ese mismo resultado. El costo es estado adicional, política de expiración y coordinación con el almacenamiento.

La evidencia apropiada sería historia de lecturas/escrituras, staleness y comparación de quórums. También debe indicarse una condición que refute la elección o haga preferible una solución más simple.
\end{uaanswer}

\begin{uaexercise}
Distinguir estado físico vs estado lógico en un sistema replicado.
\end{uaexercise}
\begin{uaanswer}
La idempotencia separa intentos físicos de una operación lógica. Se usa una clave estable, se persiste el resultado junto con el efecto y los reintentos devuelven ese mismo resultado. El costo es estado adicional, política de expiración y coordinación con el almacenamiento.

La evidencia apropiada sería historia de lecturas/escrituras, staleness y comparación de quórums. También debe indicarse una condición que refute la elección o haga preferible una solución más simple.
\end{uaanswer}

\begin{uaexercise}
Explicar por qué “una sola fuente de verdad” es difícil de sostener en presencia de particiones.
\end{uaexercise}
\begin{uaanswer}
La respuesta debe situarse en replicación síncrona y asíncrona, stale reads, quórums, causalidad y modelos de consistencia. Se comienza declarando el supuesto decisivo y el comportamiento observable; luego se recorre una ejecución nominal y otra bajo réplica atrasada, escritura concurrente, partición y failover. El mecanismo elegido debe vincularse con una garantía concreta, evitando afirmar propiedades fuera de su alcance.

El argumento debe identificar la hipótesis, construir la cadena lógica o el contraejemplo y concluir exactamente qué propiedad queda demostrada. Una intuición sin secuencia verificable no es suficiente.

La evidencia apropiada sería historia de lecturas/escrituras, staleness y comparación de quórums. También debe indicarse una condición que refute la elección o haga preferible una solución más simple.
\end{uaanswer}

\begin{uaexercise}
Definir linearizabilidad y dar un ejemplo que la viole.
\end{uaexercise}
\begin{uaanswer}
Linearizabilidad exige que cada operación parezca ocurrir en un punto entre invocación y respuesta y que el orden respete tiempo real. Para validar una historia se propone un orden secuencial compatible; si ninguna ubicación de las operaciones satisface ambos requisitos, la historia no es linearizable.

La evidencia apropiada sería historia de lecturas/escrituras, staleness y comparación de quórums. También debe indicarse una condición que refute la elección o haga preferible una solución más simple.
\end{uaanswer}

\begin{uaexercise}
Definir consistencia causal con la relación happens-before.
\end{uaexercise}
\begin{uaanswer}
La respuesta debe situarse en replicación síncrona y asíncrona, stale reads, quórums, causalidad y modelos de consistencia. Se comienza declarando el supuesto decisivo y el comportamiento observable; luego se recorre una ejecución nominal y otra bajo réplica atrasada, escritura concurrente, partición y failover. El mecanismo elegido debe vincularse con una garantía concreta, evitando afirmar propiedades fuera de su alcance.

La evidencia apropiada sería historia de lecturas/escrituras, staleness y comparación de quórums. También debe indicarse una condición que refute la elección o haga preferible una solución más simple.
\end{uaanswer}

\begin{uaexercise}
Explicar por qué dos timestamps físicos de nodos distintos no alcanzan para decidir causalidad entre eventos.
\end{uaexercise}
\begin{uaanswer}
Porque los relojes locales pueden derivar, corregirse o no estar perfectamente sincronizados. Además, una marca física menor no prueba dependencia causal: un evento solo pudo influir sobre otro si hay orden local, mensaje o transitividad que conecte ambos eventos.
\end{uaanswer}

\begin{uaexercise}
Aplicar las reglas de relojes de Lamport a una secuencia simple de eventos con envío y recepción de mensajes.
\end{uaexercise}
\begin{uaanswer}
La idempotencia separa intentos físicos de una operación lógica. Se usa una clave estable, se persiste el resultado junto con el efecto y los reintentos devuelven ese mismo resultado. El costo es estado adicional, política de expiración y coordinación con el almacenamiento.

La evidencia apropiada sería historia de lecturas/escrituras, staleness y comparación de quórums. También debe indicarse una condición que refute la elección o haga preferible una solución más simple.
\end{uaanswer}

\begin{uaexercise}
Dar un ejemplo permitido por causal pero no por linearizabilidad.
\end{uaexercise}
\begin{uaanswer}
La idempotencia separa intentos físicos de una operación lógica. Se usa una clave estable, se persiste el resultado junto con el efecto y los reintentos devuelven ese mismo resultado. El costo es estado adicional, política de expiración y coordinación con el almacenamiento.

La evidencia apropiada sería historia de lecturas/escrituras, staleness y comparación de quórums. También debe indicarse una condición que refute la elección o haga preferible una solución más simple.
\end{uaanswer}

\begin{uaexercise}
Definir consistencia eventual y sus límites.
\end{uaexercise}
\begin{uaanswer}
La respuesta debe situarse en replicación síncrona y asíncrona, stale reads, quórums, causalidad y modelos de consistencia. Se comienza declarando el supuesto decisivo y el comportamiento observable; luego se recorre una ejecución nominal y otra bajo réplica atrasada, escritura concurrente, partición y failover. El mecanismo elegido debe vincularse con una garantía concreta, evitando afirmar propiedades fuera de su alcance.

La evidencia apropiada sería historia de lecturas/escrituras, staleness y comparación de quórums. También debe indicarse una condición que refute la elección o haga preferible una solución más simple.
\end{uaanswer}

\begin{uaexercise}
Comparar los tres modelos en términos de garantías observables.
\end{uaexercise}
\begin{uaanswer}
La idempotencia separa intentos físicos de una operación lógica. Se usa una clave estable, se persiste el resultado junto con el efecto y los reintentos devuelven ese mismo resultado. El costo es estado adicional, política de expiración y coordinación con el almacenamiento.

La evidencia apropiada sería historia de lecturas/escrituras, staleness y comparación de quórums. También debe indicarse una condición que refute la elección o haga preferible una solución más simple.
\end{uaanswer}

\begin{uaexercise}
Describir leader-based replication y sus variantes sync/async.
\end{uaexercise}
\begin{uaanswer}
La respuesta debe situarse en replicación síncrona y asíncrona, stale reads, quórums, causalidad y modelos de consistencia. Se comienza declarando el supuesto decisivo y el comportamiento observable; luego se recorre una ejecución nominal y otra bajo réplica atrasada, escritura concurrente, partición y failover. El mecanismo elegido debe vincularse con una garantía concreta, evitando afirmar propiedades fuera de su alcance.

La evidencia apropiada sería historia de lecturas/escrituras, staleness y comparación de quórums. También debe indicarse una condición que refute la elección o haga preferible una solución más simple.
\end{uaanswer}

\begin{uaexercise}
Explicar multi-leader y por qué introduce conflictos.
\end{uaexercise}
\begin{uaanswer}
La idempotencia separa intentos físicos de una operación lógica. Se usa una clave estable, se persiste el resultado junto con el efecto y los reintentos devuelven ese mismo resultado. El costo es estado adicional, política de expiración y coordinación con el almacenamiento.

El argumento debe identificar la hipótesis, construir la cadena lógica o el contraejemplo y concluir exactamente qué propiedad queda demostrada. Una intuición sin secuencia verificable no es suficiente.

La evidencia apropiada sería historia de lecturas/escrituras, staleness y comparación de quórums. También debe indicarse una condición que refute la elección o haga preferible una solución más simple.
\end{uaanswer}

\begin{uaexercise}
Describir leaderless replication y el rol de $N,W,R$.
\end{uaexercise}
\begin{uaanswer}
La respuesta debe situarse en replicación síncrona y asíncrona, stale reads, quórums, causalidad y modelos de consistencia. Se comienza declarando el supuesto decisivo y el comportamiento observable; luego se recorre una ejecución nominal y otra bajo réplica atrasada, escritura concurrente, partición y failover. El mecanismo elegido debe vincularse con una garantía concreta, evitando afirmar propiedades fuera de su alcance.

La evidencia apropiada sería historia de lecturas/escrituras, staleness y comparación de quórums. También debe indicarse una condición que refute la elección o haga preferible una solución más simple.
\end{uaanswer}

\begin{uaexercise}
Probar informalmente por qué $R+W>N$ implica intersección.
\end{uaexercise}
\begin{uaanswer}
La respuesta debe situarse en replicación síncrona y asíncrona, stale reads, quórums, causalidad y modelos de consistencia. Se comienza declarando el supuesto decisivo y el comportamiento observable; luego se recorre una ejecución nominal y otra bajo réplica atrasada, escritura concurrente, partición y failover. El mecanismo elegido debe vincularse con una garantía concreta, evitando afirmar propiedades fuera de su alcance.

El argumento debe identificar la hipótesis, construir la cadena lógica o el contraejemplo y concluir exactamente qué propiedad queda demostrada. Una intuición sin secuencia verificable no es suficiente.

La evidencia apropiada sería historia de lecturas/escrituras, staleness y comparación de quórums. También debe indicarse una condición que refute la elección o haga preferible una solución más simple.
\end{uaanswer}

\begin{uaexercise}
Construir un ejemplo donde $R+W\le N$ rompe lecturas frescas.
\end{uaexercise}
\begin{uaanswer}
La idempotencia separa intentos físicos de una operación lógica. Se usa una clave estable, se persiste el resultado junto con el efecto y los reintentos devuelven ese mismo resultado. El costo es estado adicional, política de expiración y coordinación con el almacenamiento.

La evidencia apropiada sería historia de lecturas/escrituras, staleness y comparación de quórums. También debe indicarse una condición que refute la elección o haga preferible una solución más simple.
\end{uaanswer}

\begin{uaexercise}
Definir conflicto concurrente y dar un ejemplo.
\end{uaexercise}
\begin{uaanswer}
La idempotencia separa intentos físicos de una operación lógica. Se usa una clave estable, se persiste el resultado junto con el efecto y los reintentos devuelven ese mismo resultado. El costo es estado adicional, política de expiración y coordinación con el almacenamiento.

La evidencia apropiada sería historia de lecturas/escrituras, staleness y comparación de quórums. También debe indicarse una condición que refute la elección o haga preferible una solución más simple.
\end{uaanswer}

\begin{uaexercise}
Explicar LWW y sus riesgos.
\end{uaexercise}
\begin{uaanswer}
La respuesta debe situarse en replicación síncrona y asíncrona, stale reads, quórums, causalidad y modelos de consistencia. Se comienza declarando el supuesto decisivo y el comportamiento observable; luego se recorre una ejecución nominal y otra bajo réplica atrasada, escritura concurrente, partición y failover. El mecanismo elegido debe vincularse con una garantía concreta, evitando afirmar propiedades fuera de su alcance.

La evidencia apropiada sería historia de lecturas/escrituras, staleness y comparación de quórums. También debe indicarse una condición que refute la elección o haga preferible una solución más simple.
\end{uaanswer}

\begin{uaexercise}
Explicar vector clocks y cómo detectan concurrencia.
\end{uaexercise}
\begin{uaanswer}
La idempotencia separa intentos físicos de una operación lógica. Se usa una clave estable, se persiste el resultado junto con el efecto y los reintentos devuelven ese mismo resultado. El costo es estado adicional, política de expiración y coordinación con el almacenamiento.

La evidencia apropiada sería historia de lecturas/escrituras, staleness y comparación de quórums. También debe indicarse una condición que refute la elección o haga preferible una solución más simple.
\end{uaanswer}

\end{document}
\iffalse
\section{Resoluciones del taller P--F--G--S--T sobre Asteria Online}

\begin{uaexercise}
\textbf{Compra de oro con pérdida de señal.} Construya un par \(P,G,F\),
incluyendo dos fallas de cliente, y explique por qué ``nunca paga dos veces''
no es una buena garantía.
\end{uaexercise}
\begin{uaanswer}
\(P\): cobro y acreditación pueden quedar en estados distintos tras una
desconexión. \(G\): el resultado reconocido por el jugador y el oro
acreditado corresponden a una operación de cobro válida. \(F\): pérdida de
señal después del cobro y reintento al recuperar la red; también timeout por
cambio de red. \(S\) puede usar una clave idempotente estable y persistir
cobro, acreditación y resultado. ``Nunca'' es una promesa absoluta; que los
reintentos no dupliquen el efecto es una propiedad de esa solución. El costo
es metadata y reconciliación con el procesador de pagos.
\end{uaanswer}

\begin{uaexercise}
\textbf{Directorio de zonas.} Compare una solución consistente y una
disponible durante un particionamiento entre datacenters.
\end{uaexercise}
\begin{uaanswer}
\(P\): coordinar altas y ocupación del directorio global. \(F\): los
datacenters siguen vivos, pero no se comunican. \(G\): el jugador recibe una
asignación válida. Una solución consistente exige quórum o bloquea el ingreso
sin confirmación; una disponible acepta decisiones locales y reconcilia luego.
La primera prioriza consistencia y pierde progreso; la segunda disponibilidad
y arriesga vistas incompatibles o asignaciones duplicadas. La evidencia es la
historia de altas, lecturas y respuestas durante la partición.
\end{uaanswer}

\begin{uaexercise}
\textbf{Evento global.} Evite aplicar dos veces un ataque reintentado y
describa un ataque de seguridad contra el mecanismo.
\end{uaexercise}
\begin{uaanswer}
\(P\): actualizar los puntos de vida compartidos y el premio final. \(F\):
respuesta perdida y reintento, o mensajes duplicados. \(G\): cada ataque
lógico válido afecta una sola vez. \(S\): identificador de ataque por sesión,
registro durable del identificador junto con la transición y validación de
versión del jefe. Un atacante puede reenviar identificadores o fabricar
ataques; el servidor debe autenticar la sesión, vincular el ataque a una
versión y ventana válidas, y aplicar límites de ritmo. El costo es coordinación
y menor disponibilidad durante una partición.
\end{uaanswer}

\begin{uaexercise}
\textbf{Recompensa única y mercado.} Compare replicación síncrona, asíncrona y
leaderless con quórums para un objeto único.
\end{uaexercise}
\begin{uaanswer}
La garantía es que una recompensa o venta confirmada no produzca dos objetos
disponibles. El líder síncrono ofrece una decisión clara, con latencia y
bloqueo sin quórum. El asíncrono ofrece disponibilidad, pero puede mostrar
estado atrasado y necesita reconciliación. Leaderless con \(R+W>N\) fuerza
intersección, aunque aún requiere versionado y resolución de conflictos. La
condición de conveniencia depende de si se acepta bloquear, reconciliar o
arriesgar lecturas stale; la evidencia son versiones y lecturas/escrituras.
\end{uaanswer}

\begin{uaexercise}
\textbf{Caída de datacenter y migración.} Construya un PFGST distinguiendo
estado confirmado y estado en vuelo.
\end{uaexercise}
\begin{uaanswer}
\(P\): recuperar una mazmorra cuyo datacenter cayó durante el golpe final.
\(F\): caída completa con eventos en vuelo. \(G\): no perder recompensa
confirmada ni otorgar una no confirmada. \(S\): checkpoints replicados,
registro de eventos, promoción de una copia y reconciliación solo de eventos
confirmados. \(T\): el daño en vuelo puede perderse o el combate reiniciarse
durante la migración; esa brecha, más el costo de replicación, debe declararse
en la Entrega 3.
\end{uaanswer}

\begin{uaexercise}
\textbf{PFGST acumulativo.} Produzca tres filas sobre los casos de uso de
Asteria, incluyendo fallas de cliente, infraestructura y seguridad.
\end{uaexercise}
\begin{uaanswer}
Una respuesta válida usa Compra de Oro (reintento e idempotencia), Directorio
(partición y elección de CAP) y Evento Global (replay y validación del
servidor). Las garantías se mantienen como intereses del jugador y las
soluciones evolucionan al agregar fallas. Debe incluirse una alternativa
descartada y su costo; por ejemplo, aceptar ingresos locales durante una
partición se descarta para objetos únicos si la reconciliación posterior no
puede preservar la garantía.
\end{uaanswer}

\begin{uaexercise}
\textbf{Garantía o promesa absoluta.} Clasifique y reformule las frases del
ejercicio sobre ``siempre'', ``nunca'' y tiempos concretos.
\end{uaexercise}
\begin{uaanswer}
(a) y (b) son promesas absolutas y deben reformularse, por ejemplo: ``una
recompensa se acredita solo cuando el combate fue resuelto según las reglas''
y ``una operación lógica de oro no produce más de un efecto confirmado''. (c)
ya expresa un interés de negocio y puede ser una garantía, siempre que se
aclare qué significa ``confirmada''. (d) es un objetivo de rendimiento, no una
garantía de negocio.

La misma garantía puede ser amenazada por una falla de cliente (reintento tras
desconexión), de infraestructura (caída o partición con réplicas divergentes)
y de seguridad (replay deliberado). La garantía permanece fija; evolucionan la
solución y el trade-off. La respuesta debe evitar convertir la garantía en
``siempre'' o ``nunca'' para ocultar los límites de la solución.
\end{uaanswer}

\fi
